% START: Metadaten ersetzen, dann die kommentierten Beispielabschnitte ändern.
% Bauen: im Repository-Hauptordner `make hausarbeit` (siehe README.md).
\documentclass[12pt,a4paper,ngerman]{scrartcl}
\usepackage[T1]{fontenc}
\usepackage[utf8]{inputenc}
\usepackage[ngerman]{babel}
\usepackage{lmodern}
\usepackage[a4paper,margin=2.5cm]{geometry}
\usepackage{setspace}
\onehalfspacing
\usepackage{amsmath,amssymb,amsthm}
\usepackage{graphicx,booktabs}
\usepackage{pgfplots}
\pgfplotsset{compat=1.18}
\usepackage[authoryear,round]{natbib}
\usepackage{csquotes}
\usepackage[hidelinks]{hyperref}
\AtBeginDocument{\renewcommand{\refname}{Literaturverzeichnis}}
% Eigene Mathematik-Umgebungen: Nummerierung erfolgt pro Abschnitt.
\newtheorem{satz}{Satz}[section]
\theoremstyle{definition}
\newtheorem{definition}[satz]{Definition}
\newtheorem{beispiel}[satz]{Beispiel}
% BEARBEITEN: wiederverwendbare Angaben; keine Namen im Dokument suchen nötig.
\newcommand{\arbeitstitel}{Quadratische Funktionen im Mathematikunterricht}
\newcommand{\personA}{Vorname Nachname 1}
\newcommand{\personB}{Vorname Nachname 2}
\newcommand{\kurs}{Mathematik · Kurs M12}
\newcommand{\einrichtung}{Schule / Hochschule}
\newcommand{\betreuung}{Name der betreuenden Lehrkraft}
\newcommand{\abgabedatum}{Abgabedatum}
\hypersetup{pdftitle={\arbeitstitel},pdfauthor={\personA, \personB}}

\begin{document}
\begin{titlepage}
  \centering
  {\Large \einrichtung\par}
  \vspace{2cm}
  {\huge\bfseries \arbeitstitel\par}
  \vspace{0.7cm}
  {\Large Gemeinsame Hausarbeit\par}
  \vfill
  % BEARBEITEN: ggf. Matrikelnummern ergänzen; Vorgaben der Einrichtung prüfen.
  \begin{tabular}{rl}
    Verfasst von: & \personA \\
                 & \personB \\
    Kurs:        & \kurs \\
    Betreuung:   & \betreuung \\
    Abgabe:      & \abgabedatum
  \end{tabular}
  \vspace{1.5cm}
\end{titlepage}

\pagenumbering{roman}
\section*{Kurzfassung}
% BEARBEITEN: etwa 100--150 Wörter zu Frage, Vorgehen und Ergebnis.
Diese Vorlage zeigt an quadratischen Funktionen grundlegende Elemente
einer wissenschaftlichen Hausarbeit: Fragestellung, mathematische Argumentation,
Unterrichtsplanung und Quellenarbeit. Die Beispiele dienen dem Ausprobieren;
sie sind durch die eigene Untersuchung und deren Ergebnisse zu ersetzen.
\tableofcontents
% OPTIONAL: erst bei mehreren Abbildungen/Tabellen sinnvoll; auskommentieren.
% \listoffigures
% \listoftables
\clearpage
\pagenumbering{arabic}

\section{Einleitung und Fragestellung}
\label{sec:einleitung}
% BEARBEITEN: Motivation, Eingrenzung, Frage und Aufbau der Arbeit erläutern.
Wie können Lernende Nullstellen quadratischer Funktionen rechnerisch bestimmen
und grafisch deuten? Die Beispielstunde verbindet eine gemeinsame Herleitung
an der Tafel mit einer anschließenden Partneraufgabe.

% FORMAT: Leerzeile beginnt einen neuen Absatz; nicht überall \\ verwenden.
\textbf{Fettdruck} hebt sparsam einen Begriff hervor; \emph{Hervorhebung}
betont eine Formulierung. Sonderzeichen schreiben: \%, \&, \_, \#.
Eine Fußnote\footnote{Fußnoten eignen sich für ergänzende Hinweise.
Quellen zu tragenden Aussagen gehören unmittelbar an die betreffende Stelle.}
zeigt eine optionale Ergänzung.

\section{Mathematische Grundlagen}
\label{sec:grundlagen}
% QUELLEN: \citet nennt den Autor im Satz; \citep setzt Autor/Jahr in Klammern.
Ein Beispiel für einen Literaturverweis im Satz: \citet{spivak2008}
behandelt Grundlagen der Analysis. Die folgenden Rechnungen sind eigene
Minimalbeispiele und kein wörtliches Zitat aus diesem Buch.

\begin{definition}[Nullstelle]
  Eine Zahl \(x_0\in\mathbb{R}\) heißt Nullstelle einer Funktion
  \(f\colon\mathbb{R}\to\mathbb{R}\), wenn \(f(x_0)=0\).
\end{definition}

% FORMEL: equation wird nummeriert; equation* oder \[...\] ohne Nummer.
\begin{equation}
  f(x)=ax^2+bx+c, \qquad a\neq 0.
  \label{eq:quadratisch}
\end{equation}
Mit \verb|\eqref{eq:quadratisch}| verweist man automatisch auf
Gleichung~\eqref{eq:quadratisch}. \verb|\label| steht hinter der Nummernquelle,
bei Abbildungen und Tabellen also hinter \verb|\caption|.

\begin{satz}[Nullstellenformel]
  Für \(\Delta=b^2-4ac\geq 0\) liegen die reellen Nullstellen von
  \eqref{eq:quadratisch} bei
  \[
    x_{1,2}=\frac{-b\pm\sqrt{b^2-4ac}}{2a}.
  \]
  Bei \(\Delta=0\) fallen beide Werte zusammen; bei \(\Delta<0\)
  gibt es keine reellen Nullstellen.
\end{satz}
\begin{proof}
  % align richtet Zeilen an & aus. \\ beginnt eine neue Rechenzeile.
  Durch quadratische Ergänzung erhält man
  \begin{align*}
    ax^2+bx+c=0
      &\iff x^2+\frac{b}{a}x=-\frac{c}{a} \\
      &\iff \left(x+\frac{b}{2a}\right)^2
              =\frac{b^2-4ac}{4a^2}.
  \end{align*}
  Für nichtnegative Diskriminante liefert Wurzelziehen die angegebenen Werte
  (das Vorzeichenpaar berücksichtigt auch negatives \(a\)). Für negative
  Diskriminante kann das Quadrat einer reellen Zahl die rechte Seite nicht erreichen.
\end{proof}

\begin{beispiel}
  Für \(f(x)=x^2-4x+3=(x-1)(x-3)\) folgt aus der Nullproduktregel
  \(x_1=1\) und \(x_2=3\). Die Probe ergibt \(f(1)=f(3)=0\).
\end{beispiel}

\subsection{Abbildungen und Tabellen}
% ABBILDUNG: LaTeX darf figure/table verschieben; [htbp] sind Platzierungswünsche.
\begin{figure}[htbp]
  \centering
  % BEARBEITEN: Ausdruck in addplot, Achsen und Beschriftung ändern.
  \begin{tikzpicture}
    \begin{axis}[width=0.75\textwidth,height=6cm,
      axis lines=middle,xlabel={$x$},ylabel={$f(x)$},
      xmin=0,xmax=4,ymin=-1.5,ymax=3.5,grid=major,
      xtick={0,1,2,3,4},ytick={-1,0,1,2,3}]
      \addplot[blue!60!black,thick,domain=0:4,samples=60]{x^2-4*x+3};
      \addplot[only marks,mark=*] coordinates {(1,0) (3,0)};
    \end{axis}
  \end{tikzpicture}
  \caption{Graph von \(f(x)=x^2-4x+3\) mit markierten Nullstellen.
    Quelle: eigene Darstellung.}
  \label{fig:parabel}
\end{figure}
Abbildung~\ref{fig:parabel} ergänzt die Rechnung. Externe Bilder können statt
der TikZ-Grafik mit diesem Befehl eingefügt werden:
\begin{quote}
  \small\verb|\includegraphics[width=0.7\textwidth]{bild.pdf}|
\end{quote}
Die Datei liegt dann neben \texttt{main.tex}; Urheber,
Quelle und gegebenenfalls Lizenz gehören in die Bildunterschrift.

\begin{table}[htbp]
  \centering
  \caption{Wertetabelle zum Beispiel; eigene Berechnung.}
  \label{tab:werte}
  \begin{tabular}{@{}rrrrr@{}}
    \toprule
    \(x\) & 0 & 1 & 2 & 3 \\
    \midrule
    \(f(x)\) & 3 & 0 & -1 & 0 \\
    \bottomrule
  \end{tabular}
\end{table}
Tabelle~\ref{tab:werte} enthält ausgewählte Funktionswerte.

\section{Didaktische Planung der Unterrichtsstunde}
% BEARBEITEN: Klassenstufe, Voraussetzungen, Lernziele, Methoden und Zeiten.
Vorausgesetzt werden der Umgang mit Termen und das Ablesen von Koordinaten.
Als Lernziel bestimmen die Lernenden Nullstellen eines faktorisierbaren
quadratischen Terms und erklären deren Bedeutung am Graphen.
\begin{enumerate}
  \item \textbf{Einstieg (5 Minuten):} Vorwissen zu Nullstellen aktivieren.
  \item \textbf{Erarbeitung (10 Minuten):} Rechenweg gemeinsam an der Tafel entwickeln.
  \item \textbf{Übung (15 Minuten):} Aufgaben in Paaren lösen und kontrollieren.
  \item \textbf{Sicherung (10 Minuten):} Lösungen vergleichen und Fehler besprechen.
  \item \textbf{Abschluss (5 Minuten):} Nullproduktregel in eigenen Worten erklären.
\end{enumerate}
Die Anregung, einen Lösungsweg zu planen und das Ergebnis rückblickend zu prüfen,
lässt sich mit der Problemlöseheuristik bei \citet{polya1945} verbinden.
% Vor Abgabe konkrete Fundstelle in eurer Ausgabe prüfen und ergänzen:
% \citep[vgl.][S.~23--25]{polya1945} (Seiten hier nur Syntax-Beispiel!).

\subsection{Zusammenarbeit zu zweit}
% BEARBEITEN: tatsächliche Beiträge nennen; beide lesen die gesamte Endfassung.
\begin{itemize}
  \item \personA: Beispielzuständigkeit für Einstieg und Literaturrecherche.
  \item \personB: Beispielzuständigkeit für Herleitung und Aufgabenentwicklung.
  \item Gemeinsam: Unterrichtsplanung, Überarbeitung und Quellenkontrolle.
\end{itemize}

\section{Wissenschaftliches Zitieren: Minimalbeispiele}
% Nicht einfach ein Literaturverzeichnis anhängen: Aussagen im Text belegen!
Ein Verweis in Klammern lautet \citep{polya1945}; mehrere Quellen zusammen
werden mit \verb|\citep{polya1945,spivak2008}| gesetzt
\citep{polya1945,spivak2008}. Die Quellenverwaltung ist unabhängig von der
Formulierung eures Texts.

Ein direktes Zitat steht in \enquote{deutschen Anführungszeichen} und braucht
einen überprüften Wortlaut sowie eine genaue Fundstelle. Diese Formulierung
zeigt nur die Schreibweise der Anführungszeichen, kein Literaturzitat.
Für eine Seitenangabe lautet die Syntax:
\begin{quote}
  \small\verb|\citep[S.~23]{schluessel}|
\end{quote}
Für ein indirektes Zitat mit Seitenbereich:
\begin{quote}
  \small\verb|\citep[vgl.][S.~23--25]{schluessel}|
\end{quote}
% BEARBEITEN: echten Schlüssel und tatsächlich geprüfte Seiten einsetzen.
Bei Internetquellen Autor/Organisation, Seitentitel, URL und Abrufdatum
erfassen; ein auskommentiertes Muster steht in \texttt{literatur.bib}.

\section{Fazit und Reflexion}
% BEARBEITEN: Frage aus Einleitung beantworten, Grenzen und Ausblick nennen.
Das Beispiel verbindet algebraische Rechnung, Graph und Ergebniskontrolle.
In der eigenen Arbeit ist zu reflektieren, welche Lernziele erreicht wurden
und woran dies erkennbar war. Vor einer tatsächlich gehaltenen Stunde keine
Beobachtungen oder Lernergebnisse als durchgeführt ausgeben.

% LITERATUR: Nur zitierte Einträge erscheinen automatisch. Kein manuelles Tippen.
\clearpage
\addcontentsline{toc}{section}{Literaturverzeichnis}
\bibliographystyle{plainnat}
\bibliography{literatur}

\clearpage
\appendix
\section{Aufgabe und Musterlösung}
% BEARBEITEN: ergänzendes Material; im Haupttext mit \ref{sec:anhang} verweisen.
\label{sec:anhang}
Bestimmt die Nullstellen von \(g(x)=x^2-5x+6\).
\[
  g(x)=(x-2)(x-3)=0 \iff x=2\ \text{oder}\ x=3.
\]
Die Probe ergibt \(g(2)=0\) und \(g(3)=0\).

% OPTIONAL: Erklärung nur verwenden, wenn gefordert, und dann den verbindlichen
% Wortlaut eurer Einrichtung übernehmen (z. B. auch Vorgaben zu KI-Nutzung).
% \clearpage
% \section*{Eigenständigkeitserklärung}
% [Hier den vorgeschriebenen und für euch zutreffenden Wortlaut einsetzen.]
% \vspace{1cm}
% Ort, Datum: \hrulefill
% \vspace{1cm}
% \personA: \hrulefill \qquad \personB: \hrulefill
\end{document}
